\input{sol_pre.tex}
\input{q_2025-11-23.tex}
\begin{document}
\input{sol_head.tex}

\Savol{1}
$2!+4!=2+24=26$; $6!$ dan boşlab har bir faktorial $10$ ga bölinadi. Oxirgi raqam $6$.
\Javob{A}

\Savol{2}
$4{,}(6)=\frac{14}3$, \ $\frac{14}3\cdot\frac{18}7=12$, \ $12-3=9$; \ $1{,}0(5)=1+\frac5{90}=\frac{19}{18}$.
\[9\cdot\frac{19}{18}-\frac52=\frac{19}2-\frac52=7.\]
\Javob{C}

\Savol{3}
Tarvuz $x$, qovun $y$, karam $z$ ta: \ $x+y+z=100$, \ $50x+10y+z=500$ $\Rightarrow$ $49x+9y=400$.\\
$x=1$ da $9y=351$, $y=39$ ($x\ge10$ bölolmaydi, $x=1$ dan boşqa $x\equiv1\pmod 9$ yöq). \ $z=100-1-39=60$.
\Javob{B}

\Savol{4}
Yoğ olinganda yoğsiz qism özgarmaydi: $18$ litr $4\%$ li sutda yoğsiz qism $0{,}96\cdot18=17{,}28$ l.\\
Yangi sutda yoğsiz qism $64\%$: \ $0{,}64x=17{,}28$ $\Rightarrow$ $x=27$ l.
\Javob{C}

\Savol{5}
$3^m\cdot81^n=3^{m+4n}=3^3=27$.
\Javob{B}

\Savol{6}
$3{,}25+\sqrt3=\left(\sqrt3+\frac12\right)^2$, \ $7-4\sqrt3=\left(2-\sqrt3\right)^2$:
\[\sqrt3+\frac12+2-\sqrt3=\frac52.\]
\Javob{C}

\Savol{7}
$t=\sqrt[3]{9+4\sqrt5}+\sqrt[3]{9-4\sqrt5}$: $t^3=18+3\sqrt[3]{81-80}\,t=18+3t$ $\Rightarrow$ $t=3$.\\
$u=\sqrt[3]{7+5\sqrt2}+\sqrt[3]{7-5\sqrt2}$: $u^3=14+3\sqrt[3]{49-50}\,u=14-3u$ $\Rightarrow$ $u=2$.\\
Ifoda $=t-u=1$.
\Javob{B}

\Savol{8}
Geometrik: $p^4=p\cdot k$ $\Rightarrow$ $k=p^3$. \ Arifmetik: $2k=p+p^2$ $\Rightarrow$ $2p^2-p-1=0$ $\Rightarrow$ $p=1$ ($k=p$, mos emas) yoki $p=-\frac12$.\\
$k=-\frac18$, $d=k-p=\frac38$: \ $S_{10}=\frac{2\cdot\left(-\frac12\right)+9\cdot\frac38}2\cdot10=5\cdot\frac{19}8=\frac{95}8$.
\Javob{A}

\Savol{9}
$\dfrac{b_1^2(1+q^2+q^4)}{b_1(1+q+q^2)}=b_1(1-q+q^2)=\dfrac{84}{14}=6$, \ $b_1(1+q+q^2)=14$.\\
$\dfrac{1+q+q^2}{1-q+q^2}=\dfrac73$ $\Rightarrow$ $2q^2-5q+2=0$ $\Rightarrow$ $q=\frac12$ (kamayuvçi), $b_1=8$.
\[S_6=\frac{8\left(1-\frac1{64}\right)}{1-\frac12}=\frac{63}4=15\frac34.\]
\Javob{B}

\Savol{10}
$|b-a|=a-b$, $|a-c|=a-c$, $|b-c|=c-b$:
\[\frac1b-\frac1a+\frac1c-\frac1a+\frac1b-\frac1c=\frac2b-\frac2a=\frac{2a-2b}{ab}.\]
\Javob{C}

\Savol{11}
$ab^3+a^3b=ab(a^2+b^2)=y(x^2-2y)$, \ $a^2+b^2=x^2-2y$:
\[yx^2-2y^2+x^2-2y+5y+2y^2=yx^2+x^2+3y.\]
\Javob{D}

\Savol{12}
$s=\sin2\alpha$, $c=\cos2\alpha$:
\[\frac{s^2\cdot\frac{s^2}{c^2}}{\frac{s^2-c^2}{s^2}}-\frac{c^2\cdot\frac{c^2}{s^2}}{\frac{s^2-c^2}{c^2}}=\frac{s^8-c^8}{s^2c^2(s^2-c^2)}=\frac{(s^4+c^4)(s^2+c^2)}{s^2c^2}=\frac{1-2s^2c^2}{s^2c^2}=\frac1{s^2c^2}-2.\]
$+2$ bilan: $\dfrac1{s^2c^2}=\dfrac4{\sin^24\alpha}$.
\Javob{C}

\Savol{13}
$\sin x\left(\sqrt5+2\cos x\right)=0$; \ $\cos x=-\frac{\sqrt5}2$ — yeçim yöq. \ $\sin x=0$: oraliqda $\pi$ va $2\pi$, yiğindi $3\pi$.
\Javob{A}

\Savol{14}
$2^{x+2}=-3-a$, $x\ne-5$. \ Yeçim bor: $-3-a>0$ $\Rightarrow$ $a<-3$; \ $x=-5$ bölmasligi uçun $-3-a\ne2^{-3}$, ya'ni $a\ne-3{,}125$.
\Javob{A}

\Savol{15}
$x>0$, $\ln x(x+3)\ne0$. \ Surat: $\ln\dfrac{4(2x+3)}5$, maxraj: $\ln x(x+3)$:
\[\frac{4(2x+3)}5=x^2+3x\ \Rightarrow\ 5x^2+7x-12=0\ \Rightarrow\ x=1\ \left(x=-\tfrac{12}5\text{ — mos emas}\right).\]
$x=1$ da maxraj $\ln4\ne0$.
\Javob{C}

\Savol{16}
$x\ne-7$: \ $(x+2)^2=25$ $\Rightarrow$ $x=3$ yoki $x=-7$ (çetlaştiriladi). $1$ ta ildiz.
\Javob{A}

\Savol{17}
$a=x^2-x-1$, $b=x^2-3x-2$, $a+b=2x^2-4x-3$: \ $a^3+b^3=(a+b)^3$ $\Rightarrow$ $3ab(a+b)=0$.\\
$a=0$: $x=\frac{1\pm\sqrt5}2$; \ $b=0$: $x=\frac{3\pm\sqrt{17}}2$; \ $a+b=0$: $x=\frac{2\pm\sqrt{10}}2$ — hammasi har xil: $6$ ta.
\Javob{A}

\Savol{18}
$\dfrac{x^2-5x+64}{(x-4)(x-6)}+\dfrac{10}{x-4}\le0$ $\Rightarrow$ $\dfrac{x^2+5x+4}{(x-4)(x-6)}=\dfrac{(x+1)(x+4)}{(x-4)(x-6)}\le0$.\\
Yeçim: $[-4;-1]\cup(4;6)$. Butunlar: $-4,-3,-2,-1,5$; yiğindi $-5$.
\Javob{A}

\Savol{19}
$t=\sqrt[6]{x^2+3x}\ge0$: $t^2+t-6=0$ $\Rightarrow$ $t=2$ $\Rightarrow$ $x^2+3x-64=0$, $D>0$. Viyet: köpaytma $-64$.
\Javob{C}

\Savol{20}
$f(x)=a(x-2)^2(x+3)$, \ $f(0)=12a=4$ $\Rightarrow$ $a=\frac13$.
\[a+b+c+d=f(1)=\frac13\cdot1\cdot4=\frac43.\]
\Javob{C}

\Savol{21}
$\dfrac3{(x-5)(x-1)}=\dfrac34\left(\dfrac1{x-5}-\dfrac1{x-1}\right)$:
\[\frac34\left.\ln\left|\frac{x-5}{x-1}\right|\right|_2^4=\frac34\left(\ln\frac13-\ln3\right)=-\frac32\ln3.\]
\Javob{C}

\Savol{22}
$\lg\sqrt[n]{x+\sqrt{1+x^2}}=\frac1n\lg\left(x+\sqrt{1+x^2}\right)$ — toq funksiya ($(-x+\sqrt{1+x^2})(x+\sqrt{1+x^2})=1$), $\cos(kx^n)$ — juft ($\cos$ juft). Demak $f$ toq, $f'$ juft: $f'(-2023)=2024$, javob $4048$.
\Javob{D}

\Savol{23}
Aylanalar $T$ nuqtada urinadi. $T$ markazli gomotetiya kiçik aylanani kattasiga ötkazadi, $TQ=3x$ vatar $TP=5x$ vatarga ötadi: koeffitsiyent $\frac53$.\\
Segmentlar yuzi $\left(\frac53\right)^2=\frac{25}9$ marta özgaradi: katta segment $9\cdot\frac{25}9=25$.
\[S=25-9=16.\]
\Javob{A}

\Savol{24}
$f'(x)=\frac14-\frac4{x^2}=0$ $\Rightarrow$ $x=-4$: $f(-4)=-2$. \ $f(-6)=-\frac{13}6$, $f(-2)=-\frac52$. Eng kattasi $-2$.
\Javob{C}

\Savol{25}
Muntazam uçburçak içidagi nuqtadan tomonlargaça masofalar yiğindisi balandlikka teng: $h=12$.\\
$a=\frac{2h}{\sqrt3}=8\sqrt3$, \ $S=\frac{\sqrt3}4\cdot192=48\sqrt3$.
\Javob{D}

\Savol{26}
Böyalgan soha — uçi katta kvadratning yuqori çap uçida, asosi uçinçi kvadratning öng tomoni ($3$) bölgan uçburçak; balandligi (gorizontal masofa) $9+3=12$:
\[S=\frac12\cdot3\cdot12=18.\]
\Javob{C}

\Savol{27}
$\angle BAD=\angle DAC=\gamma$, $AD=DC$ $\Rightarrow$ $\angle C=\gamma$. Bissektrisa xossasi: $\frac{AB}{AC}=\frac{BD}{DC}=\frac12$.\\
Sinuslar teoremasi: $\frac{\sin\gamma}{\sin B}=\frac12$, $B=180^\circ-3\gamma$: \ $\sin3\gamma=2\sin\gamma$ $\Rightarrow$ $3-4\sin^2\gamma=2$ $\Rightarrow$ $\gamma=30^\circ$.\\
$\alpha=180^\circ-90^\circ=90^\circ$.
\Javob{B}

\Savol{28}
Muntazam sakkizburçakning taşqi burçagi $45^\circ$, qöşni tomonlar yönalişi $45^\circ$ ga buriladi.\\
$\alpha$: yuqori tomon va undan uç tomon keyingi (pastki öng) tomon yönalişlari $3\cdot45^\circ=135^\circ$ ga farq qiladi, davomlari orasidagi burçak $180^\circ-135^\circ=45^\circ$.\\
$\beta$: çap tomonga qöşni ikki tomon ($2$ qadam farq, $90^\circ$): $\beta=180^\circ-90^\circ=90^\circ$.\\
$\alpha+\beta=135^\circ$.
\Javob{C}

\Savol{29}
$\vec a\cdot(\vec b+\vec c)=\frac12+\frac12=1$, \ $|\vec b+\vec c|^2=1+1+2\cdot\frac12=3$:
\[\cos\varphi=\frac1{1\cdot\sqrt3}=\frac1{\sqrt3}.\]
\Javob{C}

\Savol{30}
$45^\circ$ da balandlik $h$ va taşqi çizilgan aylana radiusi $R$ teng: $h=R=\frac{4\sqrt6}{\sqrt2}=4\sqrt3$.\\
Oltiburçak tomoni $R$: $S=\frac{3\sqrt3}2\cdot48=72\sqrt3$. \ $V=\frac13\cdot72\sqrt3\cdot4\sqrt3=288$.
\Javob{D}

\Savol{31}
$6$ ga karrali ikki xonali sonlardan $6$ raqami borlari: $36,\ 60,\ 66,\ 96$ — $4$ ta.
\Javob{A}

\Savol{32}
Har biri $6$ ta yöldan birini tanlaydi va $5$ km radiusli aylanadagi uçlardan biriga tuşadi. Yöllar orasidagi burçak $60^\circ k$, masofa $2\cdot5\sin(30^\circ k)$: $k=1$ da $5$, $k=2$ da $5\sqrt3$, $k=3$ da $10$. Ikkinçisining yöli birinçisinikiga nisbatan $6$ holatdan $3$ tasida ($k=2,3,4$) masofa $\ge6$:
$P=\frac36=\frac12$.
\Javob{A}

\newpage
\Savol{33}
$h_{\text{sil}}=3h_1$, $h_{\text{sil}}+h_1=30$ $\Rightarrow$ $h_1=7{,}5$, $h_{\text{sil}}=22{,}5$. Kesik konus radiuslari $R$ va $4$:
\[\pi R^2\cdot22{,}5+\frac{\pi\cdot7{,}5}3\left(R^2+4R+16\right)=1000\pi\ \Rightarrow\ 25R^2+10R-960=0\ \Rightarrow\ 5R^2+2R-192=0\ \Rightarrow\ R=6.\]

\textbf{33.} $h_1=7{,}5$ — \textbf{B}.

\textbf{34.} $R=6$ — \textbf{D}.

\textbf{35.} $\pi\cdot36\cdot h_2=432\pi$ $\Rightarrow$ $h_2=12$ — \textbf{F}.
\Javob{33 — B, \ 34 — D, \ 35 — F}

\newpage
\Savol{36}
$x=0$ ildiz emas. $x^2$ ga bölamiz:
\[x^2+\frac{16}{x^2}+k\left(x+\frac4x\right)+1=0.\]
$t=x+\frac4x$, $x^2+\frac{16}{x^2}=t^2-8$: \qquad $t^2+kt-7=0$.\\
$x^2-tx+4=0$: $|t|>4$ da ikkita $x$, $t=\pm4$ da bitta ($x=\pm2$), $|t|<4$ da $x$ yöq.

a) $k=1$: $t^2+t-7=0$, $t=\frac{-1\pm\sqrt{29}}2\approx2{,}19;\ -3{,}19$. Ikkalasi ham $|t|<4$: haqiqiy ildiz yöq — $0$.

b) Yagona ildiz: bitta $t=\pm4$, ikkinçisi $|t|<4$ bölişi kerak ($t_1t_2=-7$).\\
$t=4$: $16+4k-7=0$, $k=-\frac94$, ikkinçi ildiz $-\frac74$ — mos.\\
$t=-4$: $16-4k-7=0$, $k=\frac94$, ikkinçi ildiz $\frac74$ — mos.\\
$k$ ning $2$ ta qiymati.
\Javob{a) 0; \quad b) 2}

\newpage
\Savol{37}
$u=\frac{3x}2$: \ $\tg2u+\ctg u=\dfrac{\sin2u\sin u+\cos2u\cos u}{\cos2u\sin u}=\dfrac{\cos u}{\cos2u\sin u}$.
\[\frac{\cos u}{\cos2u\sin u}=\cos^2u.\]
1) $\cos u=0$: $\sin u=\pm1$, $\cos2u=-1$ — aniqlangan, tenglik bajariladi.\\
2) $\cos u\ne0$: $\cos u\sin u\cos2u=1$ $\Rightarrow$ $\frac14\sin4u=1$ — yeçim yöq.\\
$u=\frac\pi2+\pi n$ $\Rightarrow$ $x=\frac\pi3+\frac{2\pi n}3$.

a) $[0;\pi]$: $\frac\pi3$, $\pi$ — $2$ ta.

b) Eng katta manfiy $-\frac\pi3$, eng kiçik musbat $\frac\pi3$: yiğindi $0$.
\Javob{a) 2; \quad b) 0}

\newpage
\Savol{38}
Rasmda $A$ nuqta $f$ ning ösuvçi, $B$ esa kamayuvçi qismida: $f(2)=2-a+b=5$, $f(8)=a-8+b=3$ $\Rightarrow$ $a=4$, $b=7$.\\
$A$ nuqta $g$ ning kamayuvçi, $B$ — ösuvçi qismida: $g(2)=c-2+d=5$, $g(8)=8-c+d=3$ $\Rightarrow$ $c=6$, $d=1$.

a) $\dfrac{b+d}{c-a}=\dfrac{7+1}{6-4}=4$.

b) $C=g(0)=6+1=7$, \ $D=f(0)=-4+7=3$: \ $CD=7-3=4$.
\Javob{a) 4; \quad b) 4}

\newpage
\Savol{39}
Rasmdan: $\varphi(x)>0$ — $(-8;-1)\cup(1;5)$, $\varphi(x)<0$ — $(-1;1)$; \ $\varphi'(x)>0$ — $(-8;-4)\cup(0;3)$, $\varphi'(x)<0$ — $(-4;0)\cup(3;5)$.

a) $\dfrac{x\varphi(x)}{\varphi'(x)}<0$:\\
$(-8;-4)$: $x<0$, $\varphi>0$, $\varphi'>0$ — manfiy: $-7,-6,-5$.\\
$(-4;-1)$ — musbat. \ $(-1;0)$ va $(0;1)$ — manfiy, lekin butun son yöq ($x=\pm1$ da $\varphi=0$).\\
$(1;3)$ — musbat. \ $(3;5)$: $x>0$, $\varphi>0$, $\varphi'<0$ — manfiy: $4$.\\
Jami $4$ ta.

b) $\varphi>0$ va $\varphi'>0$: $(-8;-4)\cup(1;3)$: \ $-7,-6,-5,2$ — $4$ ta.
\Javob{a) 4; \quad b) 4}

\newpage
\Savol{40}
$f$ ning grafigi $O$ dan ötadi va $Ox$ ga $x=2$ da urinadi: $f(x)=x(x-2)^2=x^3-4x^2+4x$, \ $a=-4$, $b=4$, $c=0$.\\
$g$ ning grafigi $Ox$ ga $x=0$ da urinadi va $x=2$ dan ötadi: $g(x)=x^2(x-2)=x^3-2x^2$, \ $p=-2$, $q=d=0$.

a) $\dfrac{a+b+c}{p+q+d}=\dfrac0{-2}=0$.

b) \[S=\int_0^2\bigl(f(x)-g(x)\bigr)dx=\int_0^2\left(4x-2x^2\right)dx=8-\frac{16}3=\frac83.\]
\Javob{a) 0; \quad b) $\frac83$}

\newpage
\Savol{41}
Urinmalar xossasi: $AL=\dfrac{AB+AC-BC}2$; \ $AM=\dfrac{AC}2$.

a) $LM=AM-AL=\dfrac{BC-AB}2=\dfrac{6-4}2=1$.

b) $AL=\dfrac{6+12-10}2=4$. \ $BH^2=AB^2-AH^2=BC^2-(AC-AH)^2$:
\[36-AH^2=100-(12-AH)^2\ \Rightarrow\ 24AH=80\ \Rightarrow\ AH=\frac{10}3.\]
$HL=AL-AH=4-\frac{10}3=\frac23$.
\Javob{a) 1; \quad b) $\frac23$}

\newpage
\Savol{42}
\fig{\begin{tikzpicture}[scale=1.2,font=\small]
\coordinate (D) at (0,0); \coordinate (C) at (2.798,0); \coordinate (B) at (2.798,2.798); \coordinate (A) at (0,2.798); \coordinate (E) at (0.863,0.505); \coordinate (E1) at (2.293,-1.935);
\draw (A)--(B)--(C)--(D)--cycle; \draw (D)--(E)--(B); \draw (E)--(C);
\draw[axs] (C)--(E1)--(D); \draw[ax] (E)--(E1);
\node[above left] at (A) {$A$}; \node[above right] at (B) {$B$}; \node[right] at (C) {$C$}; \node[left] at (D) {$D$}; \node[above] at (0.95,0.55) {$E$}; \node[below] at (E1) {$E'$};
\node[aux,right] at ($(C)!0.5!(E1)$) {$2$}; \node[aux,left] at ($(D)!0.5!(E1)$) {$3$};
\end{tikzpicture}}
$\triangle CEB$ ni $C$ atrofida $90^\circ$ ga buramiz: $B\to D$, $E\to E'$. \ $CE'=CE=2$, $DE'=BE=3$, $\angle ECE'=90^\circ$.

a) $\triangle ECE'$ — teng yonli töğri burçakli: $EE'=2\sqrt2$, $\angle CEE'=45^\circ$.\\
$\triangle DEE'$: $DE^2+EE'^2=1+8=9=DE'^2$ $\Rightarrow$ $\angle DEE'=90^\circ$.\\
$\angle DEC=90^\circ+45^\circ=135^\circ$.

b) $\triangle DEC$: $DC^2=1+4-2\cdot1\cdot2\cos135^\circ=5+2\sqrt2$.\\
Burişda $S_{BEC}=S_{DE'C}$, şuning uçun
\[S_{DEC}+S_{BEC}=S_{DECE'}=S_{ECE'}+S_{DEE'}=\frac12\cdot2\cdot2+\frac12\cdot1\cdot2\sqrt2=2+\sqrt2.\]
\[S_{DEB}=S_{BCD}-\left(S_{DEC}+S_{BEC}\right)=\frac{5+2\sqrt2}2-2-\sqrt2=\frac12.\]
\Javob{a) $135^\circ$; \quad b) $\frac12$}

\newpage
\Savol{43}
$\angle ADB=30^\circ$ — $AB$ yoyga tiralgan içki burçak: $\smile AB=60^\circ$. Teng vatarlar teng yoylarga tiraladi: $\smile AB=\smile BC=\smile CD=\smile DE=60^\circ$, $\smile EA=120^\circ$.\\
$AB=2R\sin30^\circ=R$ $\Rightarrow$ $R=4$. \ $BD$ va $AE$ — $120^\circ$ li yoylarga tiralgan vatarlar: $BD=AE=4\sqrt3$, ular simmetrik va parallel.\\
Demak $ABDE$ — aylanaga içki çizilgan parallelogramm, ya'ni töğri törtburçak ($4\times4\sqrt3$), $F$ — uning markazi; $S_{ABDE}=16\sqrt3$.

a) $S_{BCD}=\frac12\cdot4\cdot4\sin120^\circ=4\sqrt3$, \ $S_{BFD}=\frac14S_{ABDE}=4\sqrt3$:
\[S_{BCDF}=8\sqrt3.\]

b) $S_{AEF}=\frac14S_{ABDE}=4\sqrt3$.
\Javob{a) $8\sqrt3$; \quad b) $4\sqrt3$}

\newpage
\Savol{44}
$x+y=160$, $x-y=130$ $\Rightarrow$ $x=145$, $y=15$. Radius $r=5$ mm.

a) Konus: $V=\frac13\pi r^2y=\frac13\pi\cdot25\cdot15=125\pi$ mm$^3$.

b) Silindr: $V=\pi r^2x=\pi\cdot25\cdot145=3625\pi$ mm$^3$.
\Javob{a) $125\pi$; \quad b) $3625\pi$}

\newpage
\Savol{45}
\fig{\begin{tikzpicture}[scale=0.9,font=\small]
\coordinate (P) at (0,0); \coordinate (Q) at (5,0); \coordinate (R) at (3.2,2.4); \coordinate (H) at (3.2,0);
\fill[okgreen!12] (P)--(Q)--(R)--cycle; \draw (P)--(Q)--(R)--cycle; \draw[axs] (R)--(H);
\draw (3.2,0.2)--(3.4,0.2)--(3.4,0);
\node[below] at (1.6,0) {$3{,}2$}; \node[below] at (4.1,0) {$1{,}8$}; \node[aux,left] at (3.2,1.1) {$h$};
\node[above left] at (1.6,1.2) {$4$}; \node[above right] at (4.1,1.2) {$3$};
\end{tikzpicture}}
Tomning yon (uçburçak) yoği — töğri burçakli uçburçak, töğri burçak uçidan tuşirilgan balandlik gipotenuzani $3{,}2$ va $1{,}8$ ga böladi: gipotenuza $5$ (bu parallelepipedning ikkinçi ölçami).
\[\text{katetlar: }\sqrt{5\cdot3{,}2}=4,\qquad\sqrt{5\cdot1{,}8}=3.\]
Tunuka qism — ölçamlari $12\times4$ va $12\times3$ bölgan ikkita töğri törtburçak.

a) $S_{\text{tunuka}}=12\cdot4+12\cdot3=84$ m$^2$.

b) Parallelepiped $12\times5\times12$, $4$ ta yon yoği: $2\cdot12\cdot12+2\cdot5\cdot12=408$ m$^2$.
\[84\cdot180\,000+408\cdot80\,000=15\,120\,000+32\,640\,000=47\,760\,000\text{ söm}.\]
\Javob{a) 84; \quad b) $47\,760\,000$}

\ShowJavoblar
\end{document}
